\documentclass[11pt]{article}
\usepackage{microtype,amsmath,amssymb,amsthm,mathtools,booktabs,array}
\usepackage[margin=1.05in]{geometry}
\usepackage[colorlinks=true,linkcolor=blue!45!black,citecolor=blue!45!black,urlcolor=blue!45!black]{hyperref}
\usepackage{xcolor}
\hypersetup{pdftitle={Constant quadratic fibres:base-marker operator certificates},pdfsubject={Working mathematical manuscript}}
\newtheorem{theorem}{Theorem}[section]
\newtheorem{proposition}[theorem]{Proposition}
\newtheorem{lemma}[theorem]{Lemma}
\newtheorem{corollary}[theorem]{Corollary}
\theoremstyle{definition}\newtheorem{definition}[theorem]{Definition}
\theoremstyle{remark}\newtheorem{remark}[theorem]{Remark}
\newcommand{\E}{\mathbb E}
\newcommand{\C}{\mathbb C}
\newcommand{\tr}{\operatorname{tr}}
\newcommand{\wt}{\operatorname{wt}}
\newcommand{\G}{\mathsf G}
\DeclareMathOperator{\Span}{span}
\title{Constant quadratic fibres:\\base-marker operator certificates}
\author{}\date{17 November 2026: internal proof note}
\begin{document}\maketitle

\section{A reduction of the original markers}

Let $U=(U_1,\ldots,U_r)$ and $Y$ be independent standard Gaussian
coordinates and let
\[
 P=aY^2+B(U)Y+g(U),\qquad a\text{ constant}.
\]
The degrees of $B,g$ and the number of base variables are arbitrary.
For a polynomial marker $Q$ use the common normalization
\[
 \Phi_Q(t)=\sum_{m\ge0}\E[QP^m]\frac{t^m}{m!}.
\]
All statements concern the original Gaussian functional. No convergence
of its exponential integral is assumed.

\begin{theorem}\label{thm:actual}
Over every characteristic-zero field, the full actual polynomial-Weyl
image is generated by the periods $\Phi_Q$ with $Q\in K[U]$.
In particular $P$ is safe if and only if all these base-marker periods
are polynomial. Purity is not needed for this equivalence.

Every marker involving $Y$ has an explicit finite expression in
parameter derivatives of base-marker periods, with polynomial time
coefficients. The construction retains spatial integration-by-parts
witnesses and uses no actual image-kernel computation.
\end{theorem}

This reduces the markers, not the phase or its number of variables.
For two base variables it does not prove that a pure phase is safe.
It says that any failure already has a marker independent of the
quadratic fibre. The remaining base periods involve a genuine
parameter-dependent binary exponential, to which the fixed-phase
binary theorem cannot simply be applied.

We prove a stronger coefficient-ring statement, which identifies
the specialization issue precisely. Let $R$ be any commutative
$\mathbb Q$-algebra, without reducedness or Noetherian assumptions.
Put
\begin{align*}
 A&=R[t,U,Y],\\
 D_i&=\partial_{U_i}+t(B_iY+g_i)-U_i,
 &D_Y&=\partial_Y-(1-2at)Y+tB,\\
 M&=A\big/\bigl(\sum_iD_iA+D_YA\bigr),
 &\nabla&=\partial_t+P.
\end{align*}
This is a vector-space or $R[t]$-module quotient, not a quotient
by the ordinary Jacobian ideal. The operators $D_i,D_Y$ commute
with $\nabla$, so it is a module over $D_R=R[t]\langle\partial_t\rangle$.
Write $[Q]$ for a class, let $N$ be the $D_R$-submodule generated by
all $[Q]$ with $Q\in R[U]$, and let $I_B\subset R$ be the ideal
of all coefficients of $B$.

\begin{theorem}[Ambient certificates and coefficient content]
\label{thm:content}
Every even fibre class belongs to $N$, and
\[
 I_B(M/N)=0.
\]
These statements have universal polynomial operator certificates.
If $I_B=R$, then $M=N$ over the ring $R$ itself. If $B=0$
identically, the actual Gaussian image is still generated by base
markers, because every odd fibre period is zero by parity.
An ambient odd point-torsion module may remain when $B=0$;
it must not be mistaken for an actual Gaussian image.
\end{theorem}

Over a field, a nonzero $B$ has $I_B=K$, proving
Theorem~\ref{thm:actual}. Over a general ring, the conclusion $M=N$
requires the stated content condition. We do not infer it from
fieldwise calculations when $I_B$ is nonunit.

\section{Two polynomial certificates for every even marker}

Set
\[
 s=1-2at,\qquad u=1-at,\qquad 2u-s=1,
 \qquad E_j(Q)=[QY^{2j}].
\]
Here $Q\in R[t,U]$ may depend polynomially on time.
Fibre integration by parts gives, for every $m\ge0$,
\begin{equation}\label{eq:conditional}
 s^m[QY^m]=C_m(Q):=
 \sum_{k=0}^{\lfloor m/2\rfloor}
 \frac{m!}{2^k k!(m-2k)!}\,
       t^{m-2k}s^k[QB^{m-2k}].
\end{equation}
This is a polynomial identity in $M$; no inverse of $s$ is used.
To retain its exact spatial witness, define $W_0=0$, $W_1(Q)=-Q$
and
\[
 W_m(Q)=tW_{m-1}(BQ)+(m-1)sW_{m-2}(Q)-s^{m-1}QY^{m-1}.
\]
Then $s^mQY^m-C_m(Q)=D_YW_m(Q)$, with brackets removed
from the polynomial expression defining $C_m$.

The second identity uses the parameter connection:
\begin{equation}\label{eq:evenrec}
 uE_{j+1}(Q)=t\nabla E_j(Q)+(2j+1)E_j(Q)
                        -tE_j(gQ+Q_t).
\end{equation}
Its difference, left minus right at the polynomial level, is
$-D_Y(QY^{2j+1})$. Thus it too is an ambient identity.

Let $p_j=2j-1$ for $j\ge1$. Construct an expression $A_j(Q)\in N$
for $u^{p_j}E_j(Q)$ as follows:
\begin{align*}
 A_1(Q)&=t\nabla[Q]+[Q]-t[gQ+Q_t],\\
 A_{j+1}(Q)&=tu\nabla A_j(Q)
       +\bigl(tp_ja+(2j+1)u\bigr)A_j(Q)
       -tu A_j(gQ+Q_t).
\end{align*}
The product rule and \eqref{eq:evenrec} prove the assertion
inductively. If $V_j(Q)$ is the witness for
$u^{p_j}QY^{2j}-A_j(Q)=D_YV_j(Q)$, it is retained by
\begin{align*}
 V_1(Q)&=-QY,\\
 V_{j+1}(Q)&=tu\nabla V_j(Q)
       +\bigl(tp_ja+(2j+1)u\bigr)V_j(Q)
       -tu V_j(gQ+Q_t)-u^{p_j+1}QY^{2j+1}.
\end{align*}

Choose polynomials $\alpha_j,\beta_j\in\mathbb Z[a,t]$ with
\[
 \alpha_j u^{2j-1}+\beta_j s^{2j}=1.
\]
They exist universally, including $a=0$ and nilpotent $a$:
expand $(2u-s)^{4j-2}=1$, assigning each term with $u$ exponent
at least $2j-1$ to the first summand and every remaining term
to the second. Consequently
\begin{equation}\label{eq:evenpacket}
 E_j(Q)=\alpha_j A_j(Q)+\beta_j C_{2j}(Q)\in N.
\end{equation}
Its spatial witness is $\alpha_jV_j+\beta_jW_{2j}$.
There is no time-pole guard and no coprimality assumption on a
coefficient specialization: $2u-s=1$ is universal.

\section{Odd markers and the time commutator}

Write
\[
 R_j(Q)=[QY^{2j+1}],\qquad
 L_iQ=(U_i-tg_i)Q-\partial_{U_i}Q.
\]
The two relevant spatial identities are
\begin{align}
 tR_j(BQ)&=sE_{j+1}(Q)-(2j+1)E_j(Q),\label{eq:oddstart}\\
 R_j(L_iQ)&=tE_{j+1}(B_iQ).\label{eq:oddstein}
\end{align}
Define $K_{j,F}(Q)$ to be a base-marker expression for
$tR_j(FQ)$. Equation~\eqref{eq:oddstart} starts it at $F=B$.
If it has been constructed for $F$, then
\begin{equation}\label{eq:source-derivative}
 K_{j,F_i}(Q)=K_{j,F}(L_iQ)
                          -t^2E_{j+1}(B_iFQ)
\end{equation}
constructs it for $F_i$. Indeed
$L_i(FQ)=FL_iQ-F_iQ$, and \eqref{eq:oddstein} supplies
the other term. The new identity uses the old witness at $L_iQ$,
minus the even witness times $t^2$, and the additional exact term
$tD_i(FQY^{2j+1})$. Hence all derivatives of $B$ are covered
without division.

Let $b_\alpha$ be the coefficient of $U^\alpha$ in $B$.
Finite Taylor translation gives the identity
\begin{equation}\label{eq:coefficient}
 b_\alpha=\frac1{\alpha!}
    \sum_{\gamma\ge0}\frac{(-U)^\gamma}{\gamma!}
                         \partial^{\alpha+\gamma}B(U).
\end{equation}
Apply each derivative packet to $(-U)^\gamma Q/\alpha!\gamma!$
and sum. This produces a base-marker packet $T_{j,\alpha}(Q)$
for $b_\alpha tR_j(Q)$ with explicit spatial witnesses.
The sum is finite and all factorials are units in every
$\mathbb Q$-algebra. Multiplication by the Taylor monomial
occurs in the amplitude before applying the packet; one must
not treat multiplication by $U_i$ as an action on the top quotient.

Finally remove the factor $t$ by the Weyl commutator, retaining
the actual parameter connection:
\begin{equation}\label{eq:oddpacket}
\begin{split}
 b_\alpha R_j(Q)={}&\nabla T_{j,\alpha}(Q)
          -T_{j,\alpha}(Q_t+gQ)-aT_{j+1,\alpha}(Q)\\
          &-b_\alpha tE_{j+1}(BQ).
\end{split}
\end{equation}
To verify it expand $\nabla(tR_j(Q))$ and use
\[
 \nabla R_j(Q)=R_j(Q_t+gQ)+aR_{j+1}(Q)+E_{j+1}(BQ).
\]
Equivalently, $[\nabla,t]=1$ removes an origin-supported
quotient after it is shown that $t$ kills every odd generator.
There is no division by $t$ and no assumption that the ambient
source is torsion-free. Applying the same combination to the
retained spatial witnesses verifies \eqref{eq:oddpacket}
before passage to any image.

\begin{samepage}
Every coefficient $b_\alpha$ therefore annihilates every odd
class modulo $N$; the even classes already belong to $N$.
This proves $I_B(M/N)=0$. If
$\sum_\alpha c_\alpha b_\alpha=1$ in $R$, summing
\eqref{eq:oddpacket} with coefficients $c_\alpha$ proves $M=N$.
This is a global coefficient-ring statement under precisely the
unimodular content hypothesis, with nilpotents and zero divisors retained.
\par\end{samepage}

The Gaussian integration-by-parts map to $R[[t]]$ is horizontal
and annihilates every displayed spatial witness. Thus all these
identities descend to the actual image. If $B=0$, every odd
period is zero coefficientwise by parity in $Y$, while the
even packet remains valid. This completes both theorems.

\section{A numerical cutoff and retained marker data}

The packets give finite source/parameter expressions, not just
an existence theorem. Here is a sufficient degree estimate.
Suppose $I_B=R$, or work over a field, and suppose the base periods
have a common bound
\[
 \deg_t\Phi_Q\le C\deg_UQ+E\qquad(Q\in R[U]),\quad C,E\ge0.
\]
For $B\ne0$ put $b=\deg B$ and
$\delta=\max(1,\deg g,2b)$. Then every original marker obeys
\begin{equation}\label{eq:cutoff}
 \deg_t\Phi_Q\le(C\delta+2)\deg Q
                    +E+C\delta(2b+3)+b+5.
\end{equation}
This bound is deliberately sufficient. It does not identify a least slope.

Here are the degree counts. For a time-independent base amplitude
of degree $q$, the even packet $E_j$ uses base amplitudes of degree
at most $q+j\delta$, parameter-derivative order at most $j$, and
polynomial time coefficients of degree at most $4j-2$ for $j\ge1$.
For $A_j$, induction gives amplitude degree $q+j\deg g$ and
time-coefficient degree $2j-1$. The conditional packet uses
degree at most $q+2jb$. The two B\'ezout factors have degrees
at most $2j-1$ and $2j-2$, giving the stated even bounds.

Each source differentiation in \eqref{eq:source-derivative} adds
at most $\max(1,\deg g-1)\le\delta$ to amplitude degree and
one to time-coefficient degree. There are at most $b$ such steps.
Taylor translation adds at most $b$ to amplitude degree.
After \eqref{eq:oddpacket}, an odd marker $Q(U)Y^{2j+1}$ therefore
uses base amplitude degrees at most $q+(j+2b+3)\delta$,
time-coefficient degrees at most $4j+b+7$, and derivative order
at most $j+2$. Since its original degree is $q+2j+1$, these
bounds imply \eqref{eq:cutoff}; differentiating a polynomial
period cannot increase its degree. A finite coefficient-content
B\'ezout identity adds no time or amplitude degree.

When $a=0$, conditional integration gives the sharper direct bound
\[
 \deg_t\Phi_{Q(U)Y^m}\le C\deg_UQ+(1+Cb)m+E.
\]
If $B=0$, all odd periods vanish and the even construction gives
a sufficient bound with $\delta=\max(1,\deg g)$; the looser
formula \eqref{eq:cutoff} may be used with $b=0$.

For an individual marker, the retained packet lists every base
amplitude, normal-ordered parameter operator, spatial witness,
coefficient-content multiplier and specialization guard. If its
period is not polynomial, at least one base amplitude in that
finite packet has a nonpolynomial period. Thus an unsafe witness
can be projected to base markers with a degree bound, rather
than searching ever higher fibre powers.

\clearpage
\section{The exceptional torsion is real in the ambient source}

Take $R=\C$, $B=g=0$, $a=1$, and ignore unused Gaussian bases.
In the fibre source
\[
 \C[t,Y]/(\partial_Y-(1-2t)Y)\C[t,Y],
\]
the odd class $[Y]$ satisfies
\[
 (1-2t)[Y]=0,
\]
but is nonzero. To certify this, let $e$ generate the point
module $\delta_{t=1/2}$ with $(1-2t)e=0$. Send every even
class to zero and send $[Y^{2j+1}]$ to $\partial_t^j e$.
The spatial relations hold because
$(1-2t)\partial_t^j e=2j\partial_t^{j-1}e$, and the parameter
connection $\partial_t+Y^2$ respects the map. Its image contains
$e\ne0$. The actual Gaussian functional nevertheless sends
every odd class to zero by parity. This is an explicit torsion
witness, not an unidentified extra solution.

If $B$ varies with coefficients, the content-annihilator theorem
records the exact guard. On a chart where a selected coefficient
is invertible, its packet proves ambient generation. More generally
a coefficient-content B\'ezout identity suffices. At a specialization
where all such witnesses cease to be units, one must retain the
annihilated quotient rather than infer generation from generic rank.
The family $B=b$ illustrates the distinction: the nonzero-$b$
ambient theorem cannot be specialized by dividing by $b$ at $b=0$.
The actual image on that zero face is handled by its separate parity
identity. No arbitrary-ring conclusion is inferred by ignoring content torsion.

\section{What this settles and what remains}

For the open architecture $P=aY^2+B(U_1,U_2)Y+g(U_1,U_2)$,
all original markers are now controlled by base markers, with
polynomial operator identities and a sufficient degree transfer.
The unresolved assertion is precisely that purity forces those
base periods to be polynomial. In formal conditional notation
they are
\[
 \Phi_Q=(1-2at)^{-1/2}\E_U\!\left[
 Q\exp\!\left(tg+\frac{t^2B^2}{2(1-2at)}\right)\right].
\]
Wilson's theorem~\cite{Wilson} concerns powers of a single fixed
binary polynomial. The displayed parameter-dependent phase does
not satisfy that hypothesis merely because its unit series is one.
No solution of this remaining problem is claimed here.

The new reduction is algebraic Gaussian integration by parts,
normal ordering and a scalar B\'ezout identity. It imports no new
analytic stationary or rational-slope theorem. Classical
Sabbah/Maisonobe credit for those separate mechanisms is unchanged.
Earlier nonconvergent formal statements, all-ring packets,
Bernstein identities, torsion witnesses and specialization thresholds
remain sealed. This proof does not claim independent replay,
formal verification, external review, priority, or a complete
solution of the general Gaussian problem.

\begin{thebibliography}{9}
\bibitem{Wilson}
M.~Wilson, \emph{A face-isolation proof of the two-variable Gaussian
Moments Conjecture}, public preprint (2026), Theorem~1,
\href{https://arxiv.org/abs/2607.23887v1}{arXiv:2607.23887v1}.
\end{thebibliography}
\end{document}
